\documentclass[12pt,a4paper]{article}
\usepackage[utf8]{inputenc}
\usepackage{amsmath,amssymb,amsthm}
\usepackage{geometry}
\geometry{margin=1in}

\newtheorem{theorem}{Theorem}
\newtheorem{lemma}[theorem]{Lemma}
\newtheorem{corollary}[theorem]{Corollary}
\newtheorem{proposition}[theorem]{Proposition}
\theoremstyle{definition}
\newtheorem{definition}[theorem]{Definition}

\title{Problem 108: Diophantine Reciprocals I}
\author{Project Euler Solution}
\date{}

\begin{document}
\maketitle

\section{Problem Statement}

Find the least value of $n$ such that the number of distinct solutions of
\[
\frac{1}{x} + \frac{1}{y} = \frac{1}{n}
\]
exceeds 1000.

\section{Derivation}

\subsection{Algebraic Transformation}

Multiplying both sides by $xyn$:
\[
n(x + y) = xy
\]
Rearranging via Simon's Favorite Factoring Trick:
\[
xy - nx - ny + n^2 = n^2 \implies (x-n)(y-n) = n^2
\]

\subsection{Counting Solutions}

Setting $a = x - n$, $b = y - n$ with $a, b > 0$ and $ab = n^2$, each ordered factorization of $n^2$ gives a solution. The number of unordered pairs (with $x \leq y$, i.e., $a \leq b$) is:
\[
f(n) = \left\lceil \frac{d(n^2)}{2} \right\rceil
\]

\subsection{Divisor Count Formula}

If $n = \prod_{i=1}^{k} p_i^{a_i}$, then:
\[
d(n^2) = \prod_{i=1}^{k} (2a_i + 1)
\]

We need $d(n^2) \geq 2001$.

\subsection{Minimization Strategy}

\begin{theorem}
To minimize $n = \prod p_i^{a_i}$ subject to $\prod (2a_i + 1) \geq 2001$ with $a_1 \geq a_2 \geq \cdots \geq a_k \geq 1$, assign the largest exponents to the smallest primes.
\end{theorem}

\begin{proof}
If $p_i < p_j$ but $a_i < a_j$, swapping the exponents gives $n' < n$ since:
\[
\frac{n'}{n} = \frac{p_i^{a_j} p_j^{a_i}}{p_i^{a_i} p_j^{a_j}} = \left(\frac{p_i}{p_j}\right)^{a_j - a_i} < 1
\]
while the divisor count product remains unchanged.
\end{proof}

\subsection{Solution}

The optimal exponent sequence yields:
\[
n = 2^2 \cdot 3^2 \cdot 5 \cdot 7 \cdot 11 \cdot 13 = 180180
\]

Verification: $d(n^2) = 5 \cdot 5 \cdot 3 \cdot 3 \cdot 3 \cdot 3 = 2025$, so $f(n) = \lceil 2025/2 \rceil = 1013 > 1000$.

\[
\boxed{n = 180180}
\]

\section{Complexity}
\begin{itemize}
    \item \textbf{Optimized search:} Enumerate exponent sequences --- the search space is small since we only need $\sim 6$--8 prime factors.
\end{itemize}

\section{Answer}

\[
\boxed{180180}
\]

\end{document}
